% !TEX root = ../main.tex
\section{A Validation Benchmark: Corpus and Case Design}
\label{sec:benchmark}

Section~\ref{sec:limitations} concedes that the framework's claims are unmeasured. This
appendix specifies the corpus and case design that would measure them, at the level of
detail needed for someone else to build it. The sources below were checked for the
properties the design requires; the design has not been run, and no results are reported
anywhere in this paper.

\subsection{Why a real corpus rather than a synthetic one}

A synthetic corpus would let the author construct both the cases and the notion of
applicability they test, which is precisely the circularity that would make the results
worthless. The design therefore uses public regulatory sources that supply the ground
truth independently: supersession flagged by the issuing authority, version histories
published by the registry, and divergent indications recorded by two regulators.

\subsection{Sources}

\begin{center}
\small
\begin{tabular}{@{}>{\raggedright\arraybackslash}p{3.0cm}>{\raggedright\arraybackslash}p{5.4cm}>{\raggedright\arraybackslash}p{5.6cm}@{}}
\toprule
Source & Supplies & Verified property \\
\midrule
openFDA drug label API
  & Current US labels: indications, dosage, population sections, with
    \texttt{set\_id}, \texttt{version}, \texttt{effective\_time}
  & Serves \emph{only} the current version per \texttt{set\_id}; superseded versions are
    not retrievable \\
\addlinespace
DailyMed SPL services
  & Per-label version history and superseded label text
  & History endpoint returns all versions with publication dates; superseded text is
    retrievable by version \\
\addlinespace
EMA medicines register
  & EU authorised medicines: therapeutic indication, active substance, revision number,
    authorisation, withdrawal, suspension dates
  & Register is of \emph{centrally} authorised medicines, so coverage skews toward
    biologics and specialist products \\
\bottomrule
\end{tabular}
\end{center}

The division of labour between the first two is forced rather than chosen: openFDA
supplies structured current metadata and cannot supply history, so the valid-time and
supersession families depend on DailyMed.

\subsection{Case families}

\begin{enumerate}
  \item \textbf{Valid time.} A question whose correct answer differs by effective date,
        against a label with multiple versions. Ground truth is the version history.
  \item \textbf{Supersession.} A question answerable from a superseded label version,
        where the superseding version differs. Ground truth is the registry's own
        supersession record; the EU side additionally supplies withdrawal, revocation,
        and suspension dates.
  \item \textbf{Transaction versus valid time.} Cases where the SPL effective time and
        the publication date differ, so that a corpus current to one is not current to
        the other. This family exists because the two are separate recorded fields
        rather than a distinction the author imposed.
  \item \textbf{Population.} A question about a paediatric case answered from adult
        dosing or indication text. Ground truth is the label's own age-based
        distinction.
  \item \textbf{Indication applicability.} A question about an indication the label does
        not carry --- off-label use. Authentic, authoritative, correctly citable
        evidence that does not govern the case, which is
        \textsc{Misapplied} in the sense of Section~\ref{sec:cell-v}.
  \item \textbf{Jurisdiction.} A question answered from a US label where the EU
        authorisation for the same active substance carries a different indication, or
        the reverse. Neither source governs the other's territory.
\end{enumerate}

\paragraph{A seventh family that may not be sourceable.} Conflicting authority of equal
standing requires two sources that diverge and that the institution's precedence order
does not separate. Same-substance labels from different manufacturers do diverge
substantially, but the divergence appears to track prescription versus
over-the-counter status --- a distinction precedence \emph{can} resolve, which makes it
the resolved case rather than the contested one. Whether genuine equal-standing conflict
is available within a single regulatory class is an open question, and if the data does
not supply it the family should be reported as unsourceable rather than constructed.

\subsection{Configurations and measures}

Three configurations over the same corpus: generation without retrieval; retrieval by
semantic similarity alone; and governed retrieval with applicability filtering, coverage
checking, and the defeater search of Section~\ref{sec:governed-rag}. Measures as
specified in Section~\ref{sec:evaluation}: accuracy, citation faithfulness, the
ungrounded assertion rate decomposed by cause, the misapplied rate, abstention rate with
coverage, form divergence, and repeatability. The prediction the framework makes is
specific and falsifiable: the three configurations should separate sharply on the
ungrounded and misapplied rates while separating much less, or not at all, on accuracy
and citation faithfulness.

\subsection{Reuse conditions}

US federal sources are public domain as works of the United States government. The EMA
download page states no licence or reuse terms, and that should be resolved before
publishing derived data. Both APIs are rate-limited; the fetchers written for this design
cache every response and throttle below the documented limits.
